\documentclass[11pt]{article}
\usepackage[T1]{fontenc}
\usepackage[a4paper,margin=1in]{geometry}
\usepackage{amsmath,amssymb,booktabs}
% Number displayed equations and floats within each main section.
\numberwithin{equation}{section}
\numberwithin{figure}{section}
\numberwithin{table}{section}
% Keep the contents compact; subsections remain numbered in the text.
\setcounter{tocdepth}{1}
\setcounter{secnumdepth}{2}
\usepackage[hidelinks]{hyperref}
\setlength{\emergencystretch}{3em}
\newcommand{\code}[1]{\texttt{\detokenize{#1}}}
\title{Face-Block hp--AMG Solvers for HDG Trace Systems:\\Polynomial, Geometric, and Algebraic Hierarchies}
\author{HYBRIDGE algorithm notes}
\date{}
\begin{document}
\maketitle

The hp--AMG family combines polynomial reduction, geometric or agglomerated
h-coarsening, and algebraic multigrid on the condensed HDG trace system.
These variants share equations, modal coordinates, block storage, and a
multigrid correction framework. Their coarse spaces, transfers, smoothers,
and outer iterations differ. The family name is not a new solver selector;
not every variant requires a geometric mesh hierarchy.

The implemented native variant uses p-multigrid on a fixed face graph and
scalar classical AMG at degree zero. The historical names hp-BSR and
FB-HP-MG-PCG refer to this path; \code{fb-hp-mg-pcg} remains its selector
and flexible PCGF its current outer iteration. It requires an SPD reduced
trace operator and an appropriate SPD preconditioner. BSR denotes storage,
not a symmetry property or a guarantee of convergence.

Section~\ref{sec:shared} establishes the shared HDG and multigrid formalism.
Section~\ref{sec:spd} describes the implemented SPD solver and baseline policies.
Section~\ref{sec:nonsymmetric} proposes independent algebraic transfers,
nonsymmetric relaxation, and flexible GMRES.
Section~\ref{sec:hpchoices} develops the Poisson h/p alternatives: harmonic
agglomeration, energy-minimizing AMG, and auxiliary coarse spaces.
The latter two sections are designs, not implemented solver selectors or
expansions of backend support. The production backend and capability
reference remain authoritative. The companion
\href{hierarchies.md}{hierarchy choices and literature review} compares the
Poisson candidates and records the limits of the cited results.

\tableofcontents

\section{Shared HDG formalism}\label{sec:shared}
\subsection{Condensed trace system}\label{sec:condensation}
For element $K$, $\underline{u}^K$ contains the degrees of freedom of
$u_h|_K$, $\underline{q}^K$ contains those of the discrete flux on $K$
(all spatial components), and $\underline{f}^K$ is the element source load.
For one globally numbered face $\Gamma$, $\underline{\widehat{u}}^\Gamma$
contains the trace degrees of freedom on that face in its global orientation.
The element-boundary vector $\underline{\widehat{u}}^K$ stacks the local
face vectors $\underline{\widehat{u}}^{K,e}$ over $e\subset\partial K$.
If local face $(K,e)$ corresponds to global face $\Gamma=\Gamma(K,e)$,
its orientation/basis map $O_{K,e}$ gives
\[
 \underline{\widehat{u}}^{K,e}
 =O_{K,e}\underline{\widehat{u}}^{\Gamma(K,e)}.
\]
Thus $K$, $(K,e)$, and $\Gamma$ identify one element, one local face,
and one global face, respectively.

Vectors assembled over the mesh carry no spatial superscript. In particular,
let $\mathcal F_{\mathrm{free}}$ be the set of global faces with free trace
unknowns, and stack their coefficients in the solver's face order:
\[
 \underline{\widehat{u}}
 =\operatorname{col}_{\Gamma\in\mathcal F_{\mathrm{free}}}
   \bigl(\underline{\widehat{u}}^\Gamma\bigr).
\]
Here $\underline{\widehat{u}}$ is the assembled free trace vector,
whereas $\underline{\widehat{u}}^\Gamma$ is a single face block.
Prescribed trace coefficients are kept in the boundary data below.
The assembled load $\underline{b}$ and trace corrections likewise carry no
spatial superscript; subscripts may still label coordinates, levels, or
iterations.

Write the mixed element equations and their transmission contribution as
\[
 L_K
 \begin{bmatrix}\underline{q}^K\\\underline{u}^K\end{bmatrix}
 +B_K\underline{\widehat{u}}^K
 =\begin{bmatrix}0\\\underline{f}^K\end{bmatrix},\qquad
 C_K\begin{bmatrix}\underline{q}^K\\\underline{u}^K\end{bmatrix}
 +D_K\underline{\widehat{u}}^K=\underline{g}^K.
\]
Here $\underline{f}^K$ is the tested source load, and
$\underline{g}^K$ denotes prescribed transmission or boundary loads.
The second equation denotes the local contribution to the assembled transmission
equations; it is not an independently imposed element boundary equation.
Eliminating the flux and primal degrees of freedom gives
\[
 H_K=D_K-C_K L_K^{-1}B_K,\qquad
 \underline{h}^K=\underline{g}^K-C_K L_K^{-1}
 \begin{bmatrix}0\\\underline{f}^K\end{bmatrix}.
\]
Let $E_K$ gather free global trace coefficients with the required face
orientation, and let $\underline{\widehat{u}}_D^K$ contain prescribed boundary
data. Thus
\[
 \underline{\widehat{u}}^K
 =E_K\underline{\widehat{u}}+\underline{\widehat{u}}_D^K.
\]
The assembled condensed system is
\[
 A=\sum_K E_K^T H_K E_K,\qquad
 \underline{b}=\sum_K E_K^T
 (\underline{h}^K-H_K\underline{\widehat{u}}_D^K),\qquad
 A\underline{\widehat{u}}=\underline{b}.
\]
The condensed load $\underline{b}$ includes source and boundary
contributions and is distinct from the element source $\underline{f}^K$.
After the trace solve, reconstruct the HDG fields elementwise through
\[
 \begin{bmatrix}\underline{q}^K\\\underline{u}^K\end{bmatrix}
 =L_K^{-1}\left(
 \begin{bmatrix}0\\\underline{f}^K\end{bmatrix}
 -B_K\underline{\widehat{u}}^K\right).
\]
These are abstract elimination symbols; signs follow the chosen HDG flux
convention. Static condensation requires invertible local $L_K$, but neither
$L_K$ nor $A$ needs to be symmetric for these identities. The native solver
in Section~\ref{sec:spd} additionally takes the resulting SPD convention for $A$.
See the diffusion-reaction formulation for the actual element matrices.
A known global nullspace requires compatible loads and explicit constraint or
projection treatment before either nonsingular solver framework applies.
Pure Neumann Poisson, for example, has a constant nullspace unless such a
treatment removes it.

For a scalar trace at uniform degree $p$ on two-dimensional mesh edges,
each face has $p+1$ trace modes. Face-major storage gives dense $(p+1)\times(p+1)$ blocks $A_{fg}$
in BSR. Full row storage is retained; upper-triangle-only storage is not used.
For several trace components, retain all components of each modal group;
at degree zero this leaves a component block, not a scalar.
Triangular faces in three dimensions have $(p+1)(p+2)/2$ scalar modes.
These are mathematical extensions, not claims of current backend support.

\subsection{Modal coordinates and congruence}\label{sec:coordinates}
On the reference edge $\xi\in[-1,1]$, Legendre polynomials satisfy
\[
 \int_{-1}^1 P_iP_j\,d\xi=\frac{2}{2j+1}\delta_{ij},\qquad
 \psi_j=s_jP_j,\quad s_j=\sqrt{\frac{2j+1}{2}}.
\]
This is reference-edge orthonormality. It does not include a face-length
normalization; physical geometry remains in the assembled coefficients.
Use $\underline{\widehat{u}}$ for the assembled trace coefficients in $P_j$
and $\underline{\widehat{u}}_{\mathrm{m}}$ for the same assembled trace in
$\psi_j$. With $N_f=|\mathcal F_{\mathrm{free}}|$, define
$S=I_{N_f}\otimes\operatorname{diag}(s_0,\ldots,s_p)$. Then
\[
 \underline{\widehat{u}}=S\underline{\widehat{u}}_{\mathrm{m}},
 \qquad \mathcal A=SAS,\qquad
 \underline{b}_{\mathrm{m}}=S\underline{b},\qquad
 \mathcal A\underline{\widehat{u}}_{\mathrm{m}}
 =\underline{b}_{\mathrm{m}}.
\]
For SPD $A$, this congruence preserves symmetry and positive definiteness since
$v^T\mathcal A v=(Sv)^TA(Sv)>0$ for every nonzero test vector $v$.
The coordinate change itself is also valid for nonsymmetric $A$.
A warm assembly-basis trace guess maps by $S^{-1}$ into modal coordinates;
the returned modal trace maps back by $S$ before reconstructing
$\underline{u}^K$ and $\underline{q}^K$.
Reversing an edge acts diagonally by $(-1)^j$;
canonical orientation must be consistent before forming the hierarchy.

More generally, if $T$ maps normalized modal coefficients to assembly
coefficients, use
\[
 A_0=T^TAT,\quad \underline{b}_0=T^T\underline{b},
 \quad\underline{\widehat{u}}
       =T\underline{\widehat{u}}_0.
\]
This is an invertible change of trial and test coordinates for nonsymmetric
matrices too. Positive diagonal equilibration may be added with consistent
coordinate and load maps.
Arbitrary dense face changes can mix low and high modes; after such a change,
transform the polynomial embeddings instead of blindly truncating coordinates.

For the native Legendre normalization above, $T=S$ and $A_0=\mathcal A$.
The generalization in Section~\ref{sec:nonsymmetric} uses this same coordinate
interface, including when a full face-basis transform is required.

\subsection{Trace correction spaces and transfers}\label{sec:hierarchy}
Let $A_\ell$ act on $n_\ell$ trace correction coefficients at level $\ell$.
The grouping may describe modes on original faces, traces on macro-faces,
or algebraic aggregate candidates. Define
\[
 P_\ell:\mathbb R^{n_{\ell+1}}\longrightarrow\mathbb R^{n_\ell},\qquad
 R_\ell:\mathbb R^{n_\ell}\longrightarrow\mathbb R^{n_{\ell+1}},\qquad
 A_{\ell+1}=R_\ell A_\ell P_\ell.
\]
After smoothing a trace correction $e_\ell$, form
$r_c=R_\ell(r_\ell-A_\ell e_\ell)$ and update
$e_\ell\leftarrow e_\ell+P_\ell B_{\ell+1}r_c$.
The recurrence permits p-, h-, and algebraic coarsening at different levels.
Polynomial truncation changes modes on a fixed face graph; geometric
agglomeration changes the skeleton; algebraic coarsening need not correspond
to a literal coarse mesh. Fixed-block BSR remains appropriate where each
level has uniform block dimensions, which may differ between levels.

For SPD Galerkin variants, take $R_\ell=P_\ell^T$ in assembled coefficient
coordinates. Independent $R_\ell$ and $P_\ell$ give a Petrov--Galerkin variant;
transpose restriction is not assumed in Section~\ref{sec:nonsymmetric}.
Residual restriction acts on a dual load vector. In compatible nested
function spaces, field projection instead has coefficient representation
\[
 Q_{L^2}=M_c^{-1}P^TM_f,\qquad M_c=P^TM_fP,
\]
where $M_f,M_c$ are the physical mass matrices. On a straight edge of length
$L_f$, the reference-normalized modal mass is $(L_f/2)I$.
The SPD energy projection is $(P^TAP)^{-1}P^TA$.
Do not replace $P^T$ on dual residuals with a mass-weighted field projection
without reformulating the operator and load coordinates consistently.

All variants retain the original assembled equation for final acceptance.
The outer iteration must match the complete operator/preconditioner pair;
a common block layout does not transfer SPD theory to a nonsymmetric cycle.

\section{Implemented SPD pMG--AMG variant}\label{sec:spd}
\subsection{Polynomial hierarchy and multigrid cycle}
Let $p_0=p>p_1>\cdots>p_L=0$. At adjacent levels define the facewise injection
\[
 J_\ell=\begin{bmatrix}I_{p_{\ell+1}+1}\\0\end{bmatrix},\qquad
 P_\ell=I_{N_f}\otimes J_\ell,\qquad R_\ell=P_\ell^T.
\]
Prolongation fills higher modes with zeros; restriction truncates them.
With $A_0=\mathcal A$, the Galerkin operator is
\[
 A_{\ell+1}=P_\ell^T A_\ell P_\ell.
\]
Every coarse face block is therefore the leading principal subblock of its
fine block. The implementation extracts these subblocks directly from the
finest normalized BSR data, which is equivalent to recursive nested extraction.
The face graph stays fixed throughout p-coarsening. This is a restricted
high-order trace operator, not a newly assembled and condensed lower-order HDG
problem. Full-column-rank injection preserves SPD on every p-level.
At $p_L=0$ there is one scalar unknown per face; AMG then coarsens algebraically.

On each nonterminal level use $M_\ell=\operatorname{blockdiag}((A_\ell)_{ff})$.
For SPD $A_\ell$, each principal face block is SPD. The code symmetrizes stored
diagonal blocks and their computed inverses to reduce roundoff asymmetry;
this does not repair a nonsymmetric global operator.
Let $\underline{\delta\widehat{u}}_\ell$ denote the assembled modal trace
correction on level $\ell$, and let $r$ be its residual right-hand side.
A smoothing stage is
\[
 \underline{\delta\widehat{u}}_\ell
 \leftarrow\underline{\delta\widehat{u}}_\ell
 +\omega_j M_\ell^{-1}
 (r-A_\ell\underline{\delta\widehat{u}}_\ell).
\]
For polynomial order $m$ and interval $[a_\ell,b_\ell]$, the weights are
\[
 \omega_j=\left[\frac{b_\ell+a_\ell}{2}
 -\frac{b_\ell-a_\ell}{2}\cos\frac{(2j+1)\pi}{2m}\right]^{-1},
 \qquad j=0,\ldots,m-1.
\]
One sweep comprises all $m$ stages. The corresponding error polynomial is
$q_m(t)=\prod_j(1-\omega_jt)$. Setup estimates the largest eigenvalue of
$M_\ell^{-1}A_\ell$ using power iteration and the generalized Rayleigh quotient
$v^TA_\ell v/(v^TM_\ell v)$. Current defaults use 12 iterations,
$b_\ell=1.2\widehat\lambda_{\max}$, and $a_\ell=0.1b_\ell$.
These are estimates, not certified spectral bounds. The low endpoint selects
the smoothing band; it is not an assertion that all eigenvalues exceed it.
Weights are retained and fixed during repeated applications.

Denote the level-$\ell$ correction operator by $B_\ell$. Its recursive action
on $r$ is:
\begin{enumerate}
\item If $\ell=L$, return one zero-initialized, fixed-work scalar AMG cycle.
\item Set $\underline{\delta\widehat{u}}_\ell=0$ and apply the configured pre-smoothing sweeps in forward
Chebyshev stage order.
\item Form $r_c=P_\ell^T(r-A_\ell\underline{\delta\widehat{u}}_\ell)$.
\item Recursively compute
$\underline{\delta\widehat{u}}_{\ell+1}=B_{\ell+1}r_c$ and set
$\underline{\delta\widehat{u}}_\ell\leftarrow
\underline{\delta\widehat{u}}_\ell+
P_\ell\underline{\delta\widehat{u}}_{\ell+1}$.
\item Apply the matching post-smoothing sweeps in reversed stage order and
return $\underline{\delta\widehat{u}}_\ell$.
\end{enumerate}
Zero initialization is per preconditioner application; it is independent of
any warm guess for the outer solve. Fixed work avoids a tolerance-dependent
inner solve. The hierarchy and workspaces are reused while the operator is
unchanged.

At degree zero the baseline AMGX configuration uses classical AMG, PMIS
selection, AHAT strength with threshold $0.40$, D2 interpolation, L1-Jacobi
smoothing, and a dense LU terminal solve. It disables aggressive levels and
error scaling, and uses dense-LU row thresholds 128/256. The dense factorization
algorithm being LU does not itself imply a nonsymmetric input operator.
AMGX owns the scalar algebraic hierarchy; native code owns the p-levels.

\subsection{Symmetry, outer iteration, and convergence}
Let $H$ be the correction operator of the complete pre-smoothing sequence on
one level. Reversing symmetric elementary correction stages gives the adjoint
post-smoothing correction $H^T$. Writing $C=PB_cP^T$, the resulting cycle is
\[
 B=H+H^T-H^TAH+(I-H^TA)C(I-AH).
\]
Thus $B$ is symmetric if $A$ and $B_c$ are symmetric. If $B_c$ is positive
semidefinite, the last term is positive semidefinite. With $E=I-HA$,
\[
 A(H+H^T-H^TAH)A=A-E^TAE.
\]
Strict contraction of $E$ in the $A$-energy norm is a sufficient condition for
positive definiteness of the smoothing contribution and hence of $B$.
Suitable Chebyshev spectral bounds support this condition; inaccurate bounds
can invalidate it. Balanced sweeps alone are not a proof of SPD, and a fixed
AMG cycle alone is not a proof of self-adjointness of its complete action.

Production setup samples preconditioner adjointness through
\[
 \delta_B=\frac{|u^TBv-v^TBu|}
 {\|u\|_2\|Bv\|_2+\|v\|_2\|Bu\|_2}
\]
and positive action through $u^TBu/(\|u\|_2\|Bu\|_2)$, using deterministic
vectors and a tiny denominator floor. Nonfinite or unacceptable results reject
the native attempt. Iteration also checks positive preconditioned residual
and operator curvature. These samples can detect failures but cannot certify
symmetry or positive definiteness for all vectors.

The flexible CG update below does not remove the SPD operator requirement.
Section~\ref{sec:nonsymmetric} develops a proposed nonsymmetric extension by
changing the outer iteration, smoothing, and coarse transfers. That design
requires separate qualification; it does not change the current support of
this SPD implementation.

All recurrence vectors below use normalized modal coordinates. The iterate
$\underline{\widehat{u}}_{\mathrm{m},k}$ is the assembled HDG trace vector;
$r_k$, $z_k$, and $d_k$ are algebraic residual, preconditioned residual, and
trace search-direction vectors. Set
$r_0=\underline{b}_{\mathrm{m}}
-A_0\underline{\widehat{u}}_{\mathrm{m},0}$, $z_0=B_0r_0$, $d_0=z_0$, and $\rho_0=r_0^Tz_0$.
For each iteration, before any true-residual replacement,
\begin{align*}
 t_k&=A_0d_k,& \alpha_k&=\frac{\rho_k}{d_k^Tt_k},\\
 \underline{\widehat{u}}_{\mathrm{m},k+1}
 &=\underline{\widehat{u}}_{\mathrm{m},k}+\alpha_kd_k,&
 r_{k+1}&=r_k-\alpha_kt_k,\\
 z_{k+1}&=B_0r_{k+1},& \rho_{k+1}&=r_{k+1}^Tz_{k+1},\\
 \beta_k&=\frac{z_{k+1}^T(r_{k+1}-r_k)}{\rho_k},&
 d_{k+1}&=z_{k+1}+\beta_kd_k.
\end{align*}
Here $t_k$ is an operator-action work vector; the physical flux degrees of
freedom remain $\underline{q}^K$ and are reconstructed after the trace solve.
The beta formula is the implemented flexible Polak--Ribiere-style PCGF update;
it is not the ordinary $\rho_{k+1}/\rho_k$ update. Nonfinite or nonpositive
$\rho_k$ or $d_k^TA_0d_k$ causes breakdown. A residual-gap restart instead sets
$d_{k+1}=z_{k+1}$.

The acceptance contract is measured in the original assembly coefficients:
\begin{equation}\label{eq:acceptance}
 \|\underline{b}-A\underline{\widehat{u}}\|_2
 \leq\max(\mathrm{atol},\mathrm{rtol}\|\underline{b}\|_2).
\end{equation}
Since $r=S(\underline{b}-A\underline{\widehat{u}})$, a recursive modal residual must be measured as
$\|S^{-1}r\|_2$, not $\|r\|_2$. This Euclidean coefficient norm is also distinct
from a physical finite-element $L^2$ norm.
In the diffusion integration, true checks explicitly compute $\underline{b}-A(S\underline{\widehat{u}}_{\mathrm{m}})$ using
the original assembled action, measure its norm, and map it back with $S$.
This avoids acceptance based only on roundoff in the separately stored
congruence. Standalone callers may supply \code{assembly_matvec}; without it,
the check uses the stored normalized operator and the inverse scaling.

The default refresh period is ten iterations, with additional checks at
apparent convergence and exit. A true/recursive residual gap exceeding 10\%
of the true norm restarts the search direction. The solver retains one best
true-residual checkpoint and stops a stalled attempt after six true checks
without 1\% improvement. Returning that checkpoint does not relax the target.
The diffusion integration performs an independent final check. If the native
attempt fails its runtime, symmetry, curvature, or convergence gates, the
integration can construct and reuse its established hybrid AMGX fallback.
That recovery policy does not establish nonsymmetric support for this method.

\subsection{Baseline policies and implementation}
\begin{center}
\begin{tabular}{lll}
\toprule
Setting & Standard & Robust\\
\midrule
p-schedule & $p\to0$ & $p\to\lfloor p/2\rfloor\to\cdots\to0$\\
Chebyshev order & 2 & 4\\
p-level pre/post sweeps & $1+1$ & $2+2$\\
Scalar AMG pre/post sweeps & $1+1$ & $2+2$\\
AMGX coarsest sweeps & 2 & 4\\
Scalar AMG cycles per application & one V-cycle & one V-cycle\\
\bottomrule
\end{tabular}
\end{center}
For example, the robust degree-six schedule is $6\to3\to1\to0$.
Policy overrides can change the polynomial order, balanced sweep counts, and
scalar V/W cycle selection while retaining fixed work. These are baseline
configuration values, not mathematical requirements or universal optima.

Paths are relative to the repository root.
\begin{itemize}
\item \code{hybridge/linalg/multigrid/face_hp.py}:
\code{FaceBlockPmgPrototype} constructs and applies the production p-hierarchy
(despite its historical class name); \code{AmgxScalarVcycle} wraps the scalar
cycle; \code{FaceBlockHpMgPcgSolver} owns coordinate transforms, PCGF, checks,
and persistent workspaces; \code{_pcgf_beta} defines the recurrence.
\item \code{hybridge/linalg/multigrid/policy.py}: baseline and override policies.
\item \code{hybridge/linalg/gpu/legendre_face_bsr.py}:
\code{LegendreFaceBsrOperator}, \code{principal_modal_bsr_data}, and
\code{modal_degree_schedule} implement BSR actions and nested extraction.
\item \code{hybridge/solvers/diffusion_reaction.py}: assembled action,
operator reuse, field reconstruction, and native fallback integration.
\end{itemize}
The production note is \code{docs/backends/face_hp_mg_pcg.md}.
Historical measurements and design alternatives are in
\code{docs/development/plans/face_block_hp_multigrid.md}.
The HDG element derivation is
\code{docs/algorithms/diffusion_reaction/assembly.tex}.

\section{Proposed nonsymmetric hp--AMG variant}\label{sec:nonsymmetric}
Use right-preconditioned flexible GMRES (FGMRES), with a face-block
Petrov--Galerkin p-multigrid cycle and a nonsymmetric algebraic multigrid
hierarchy below the polynomial levels. Construct transfer operators from
local matrix couplings, using approximate ideal restriction/interpolation;
use nonsymmetric block or patch solves for relaxation. AIR denotes
\emph{approximate ideal restriction}. This is an algebraic alternative to
combining separate transport and diffusion preconditioners.

This section generalizes the hierarchy of Section~\ref{sec:spd}, using the condensed
system and coordinates of Section~\ref{sec:shared}. It specifies a design, not an
implemented solver selector; the existing \code{fb-hp-mg-pcg} path remains
SPD-oriented.
The proposal needs the condensed matrix, face grouping, polynomial degree,
and optional element-to-face connectivity. It does not need velocity,
diffusion tensors, a Peclet number, or a classification of the PDE.
``General'' means no SPD or particular PDE-regime assumption; it does not
mean uniform convergence for every nonsingular matrix. Local/coarse
solvability, coarse-space quality, and memory growth still matter.

There are two separable changes. FGMRES and nonsymmetric local solves remove
the SPD restriction. Algebraic Petrov--Galerkin transfers address a second
problem: polynomial truncation alone can lose essential couplings, or even
produce a singular coarse matrix from an invertible fine matrix.

Consider two trace modes ordered as retained $C$, eliminated $H$:
\[
 A=\begin{bmatrix}0&1\\-1&2\end{bmatrix},\qquad
 \det A=1,\qquad
 P_0=\begin{bmatrix}1\\0\end{bmatrix}.
\]
Despite invertibility, $P_0^TAP_0=0$. Also $(A+A^T)/2$ is singular.
An outer GMRES iteration does not fix either preconditioner setup failure.
Eliminating the high mode instead gives the nonsingular Schur complement
$0-1(2^{-1})(-1)=1/2$.

Even when truncation works, nonsymmetry removes the SPD justification for
the current Chebyshev spectral estimates and convergence argument.
Symmetrizing diagonal blocks removes information and can create singular
blocks. Replace those
steps, the positive-curvature CG recurrence, and the SPD-specific scalar
coarse policy. BSR products, modal transforms, workspaces, and residual
accounting remain useful.

\subsection{Block elimination and local transfers}
At one p-level retain all modes through $p_c$ on each face ($C$) and denote
the remaining modes by $H$. This is a split of \emph{trace modes}; the
interior $\underline{u}^K$ and $\underline{q}^K$ are already condensed.
For the derivation only, permute $H$ before $C$:
\[
 A_\ell=\begin{bmatrix}A_{HH}&A_{HC}\\A_{CH}&A_{CC}\end{bmatrix}.
\]
Assuming $A_{HH}$ is invertible, define
\begin{align*}
 P_*&=\begin{bmatrix}-A_{HH}^{-1}A_{HC}\\I\end{bmatrix},&
 R_*&=\begin{bmatrix}-A_{CH}A_{HH}^{-1}&I\end{bmatrix},\\
 \mathcal S_\ell&=A_{CC}-A_{CH}A_{HH}^{-1}A_{HC}.&&
\end{align*}
Then
\[
 A_\ell P_*=\begin{bmatrix}0\\\mathcal S_\ell\end{bmatrix},\qquad
 R_*A_\ell=\begin{bmatrix}0&\mathcal S_\ell\end{bmatrix},\qquad
 R_*A_\ell P_*=\mathcal S_\ell.
\]
Here $\mathcal S_\ell$ is the high-mode Schur complement, distinct from the
modal scaling matrix $S$ in Section~\ref{sec:coordinates}.
In general $R_*\ne P_*^T$. With $Q_H=[I\;0]^T$,
\begin{equation}\label{eq:exact}
 A_\ell^{-1}=Q_HA_{HH}^{-1}Q_H^T+P_*\mathcal S_\ell^{-1}R_*.
\end{equation}
This is the exact block-elimination limit of the proposed cycle and requires
no positivity. It provides a useful diagnostic oracle; exact global
$A_{HH}^{-1}$ is not the proposed computational method.

The retained unknowns remain the low modal coordinates. Their prolongation
now includes algebraically determined high-mode responses, potentially on
nearby faces. Thus this is a p-selected algebraic coarse space, rather than
the pure polynomial subspace of the current truncation hierarchy. Each next
level must be computed recursively from its actual parent matrix.

Choose a candidate sequence, for example $p=6\to3\to1\to0$.
It is a setup preference, not an obligation to accept a bad split.
For each coarse face group $c$, choose a bounded neighborhood $\mathcal N_c$
of high-mode face groups. A simple starting pattern is the face itself plus
neighbors in the union of the incoming and outgoing block graphs.
For an implementable first prototype, solve the small dense systems
\begin{align}
 A_{HH}[\mathcal N_c,\mathcal N_c] W[\mathcal N_c,c]
   &=-A_{HC}[\mathcal N_c,c],\label{eq:localp}\\
 Z[c,\mathcal N_c] A_{HH}[\mathcal N_c,\mathcal N_c]
   &=-A_{CH}[c,\mathcal N_c].\label{eq:localr}
\end{align}
Use pivoted LU and transposed solves, reusing a patch factor when the
neighborhoods agree. Store zero weights outside the selected pattern, set
\[
 P_\ell=\begin{bmatrix}W\\I\end{bmatrix},\qquad
 R_\ell=\begin{bmatrix}Z&I\end{bmatrix},\qquad
 A_{\ell+1}=R_\ell A_\ell P_\ell.
\]
Block sparsity decisions retain whole mode/component groups. Restriction
rows and interpolation columns may ultimately use different directed
neighborhoods. All products use the full nonsymmetric matrix.

Local restriction of this kind is motivated by $\ell$AIR
\cite{lair}; combining it with a polynomial-selected $C/H$ split here is
the proposed HYBRIDGE adaptation. Its effectiveness is not established by
the existing SPD pMG implementation.

Compute the actual product $R_\ell A_\ell P_\ell$. For example, choosing
$W=-FA_{HC}$ and $Z=-A_{CH}F$ with an approximate inverse $F$ yields
\[
 RAP=A_{CC}-2A_{CH}FA_{HC}+A_{CH}FA_{HH}FA_{HC}.
\]
It is generally \emph{not} $A_{CC}-A_{CH}FA_{HC}$. The latter defines a
different approximate block-factorization preconditioner unless additional
identities hold. The local construction need not have a single common $F$.

Unlike principal truncation, these products can add face-graph edges.
Measure transfer storage, stored block entries at all levels, setup work,
and cycle cost. Start with undropped $RAP$ so the projection identity is
testable. Any subsequent dropping needs its own accuracy and complexity
assessment. Detect singular or badly conditioned patch factors and coarse
operators; possible responses are enlarging a patch, retaining more modes,
or stopping p-coarsening and handing the current block matrix to AMG.
For instance $\left[\begin{smallmatrix}1&1\\1&0\end{smallmatrix}\right]$
is invertible but its second-mode $A_{HH}$ is zero. There is no valid exact
elimination for that prescribed split. Do not force degree zero or hide this
failure with an unexplained diagonal shift. These remedies can also fail;
report setup failure if no acceptable hierarchy fits the resource budget.

\subsection{Relaxation and algebraic coarse levels}
At a p-level relax the $H$ equations with $C$ fixed. The least expensive
choice solves each face's high-mode diagonal block without symmetrization.
A stronger option groups high modes from several connected faces.
If $G_s$ gathers a high-mode patch, factor
$A_s=G_sA_{HH}G_s^T$ and apply
\[
 B_Hr_H=\sum_s G_s^TD_s A_s^{-1}G_sr_H,\qquad
 \sum_s G_s^TD_sG_s=I.
\]
Diagonal overlap weights $D_s$ give weighted additive Schwarz; assigning
each degree of freedom to one owner gives restricted additive Schwarz.
No Cholesky factorization or positive quadratic form is required.
Larger patches may use block ILU with monitored pivots and fill.
Existing element-face and BSR patch-gather helpers can supply patch data.

Plain block Jacobi/Richardson is a cheap comparator, but need not contract
on a nonsymmetric matrix. A general starting smoother is two right-
preconditioned GMRES steps on $A_{HH}$ using $B_H$, starting at zero,
followed by a correction of the high-mode trace. This minimizes the
high-equation residual over the local Krylov search space; it guarantees
neither reduction of the full-level residual nor multigrid convergence.
It also makes the cycle nonlinear in its right-hand side, even with a
fixed number of steps, which is a reason to use outer FGMRES.
Full-face patch smoothing can be added if high-mode-only smoothing leaves
poorly represented error, with its cost recorded separately.

One p-level cycle on load $r_\ell$ is:
\begin{enumerate}
 \item Set $\underline{\delta\widehat{u}}_\ell=0$.
 \item Correct its high modes using the chosen smoother on
 $Q_H^T(r_\ell-A_\ell\underline{\delta\widehat{u}}_\ell)$.
 \item Recompute the full residual and restrict it with $R_\ell$.
 \item Apply the next-level cycle to this restricted load.
 \item Add its prolongation with $P_\ell$ to the trace correction.
 \item Recompute the residual and post-relax its high equations.
\end{enumerate}
For fixed linear high-mode correction operators $H_{\rm pre}$ and
$H_{\rm post}$, the error propagator is
\[
 E_\ell=(I-H_{\rm post}A_\ell)
 (I-P_\ell B_{\ell+1}R_\ell A_\ell)
 (I-H_{\rm pre}A_\ell).
\]
With exact high solves, ideal transfers, and an exact Schur solve, this
vanishes by \eqref{eq:exact}. Approximate components have no such automatic
guarantee. Also, $\rho(E)<1$ alone does not describe transient residual
growth for a nonnormal matrix.

The polynomial terminal matrix is the input to AMG, not necessarily the
smallest system. At this point coarsen \emph{faces}, treating all remaining
modes/components on a face as a node, and construct independent sparse
$R$ and $P$ with $A_c=RAP$. When p-coarsening reaches zero for a scalar
trace, this becomes scalar nonsymmetric AMG. Use a pivoted direct solve
only on the eventual small bottom matrix.

Block $\ell$AIR is the first reference candidate. Its local restriction
construction and block treatment have published nonsymmetric DG evidence
\cite{lair}; AIR also has direct HDG trace-system evidence in space-time
advection--diffusion \cite{hdgair}. These results support the direction,
not the claim that this particular p hierarchy is already qualified.

For a broader AMG endpoint, constrained $\ell$AIR is especially relevant:
it augments local ideal-transfer approximations with mode constraints to
improve performance and complexity across advective and diffusive cases
\cite{clair}. In a proposed block adaptation, right and left candidate
vectors should be treated separately, conceptually seeking
$PV_c\approx V$ and $R^TY_c\approx Y$. They can come from inexpensive
algebraic relaxation on $A$ and $A^T$; known near-nullspace information is
optional. The constant scalar trace, if used, occupies the degree-zero
modal coefficient on each face, not every modal entry. These algebraic
constraints do not require selecting a transport or diffusion branch.
Candidate generation, constraints, block coarsening, and complexity control
need validation together; adding a constant vector alone is not a general
convergence guarantee. This note does not claim that the local systems
\eqref{eq:localp}--\eqref{eq:localr} implement constrained $\ell$AIR.

PyAMG documents a scalar \code{air_solver} and hypre exposes AIR restriction
options \cite{pyamg,hypre}. Neither implies that this repository has a
native GPU block AIR endpoint. Current AMGX coarse configuration and any
new AIR backend must be assessed separately; changing the coarse outer
solver to FGMRES alone does not construct nonsymmetric transfer operators.

\subsection{Flexible outer iteration and convergence}
Run right-preconditioned FGMRES on the original assembled equation. A modal
cycle $\mathcal B_j$ supplies the assembly-coordinate action
\[
 z_j=T\mathcal B_j(T^Tv_j),\qquad w_j=Az_j.
\]
Arnoldi orthogonalizes $w_j$ against the residual basis $V_j$. Store the
actual correction vectors $Z_m=[z_1,\ldots,z_m]$. With
$\beta=\|\underline{b}-A\underline{\widehat{u}}_0\|_2$ at the
start of the restart cycle, solve the small least squares problem and update
\[
 y_m=\arg\min_y\|\beta e_1-\overline H_my\|_2,\qquad
 \underline{\widehat{u}}_{m}
 =\underline{\widehat{u}}_{0}+Z_my_m.
\]
This accommodates changes and inner Krylov iterations in the preconditioner
\cite{fgmres}. Do not reuse a left-preconditioned GMRES implementation as
though it were flexible right-preconditioned GMRES. In particular the
existing \code{linalg/gpu/gmres.py} interface and workspace need extension to
store $Z_m$ as well as $V_m$.

Apply the same original-system acceptance contract \eqref{eq:acceptance}
before declaring convergence, as well as at restarts and when Arnoldi
estimates become suspect. This residual uses the original load, operator,
and assembly coordinates. A small preconditioned or equilibrated residual
does not suffice. Preserve the best verified iterate and report numerical
breakdown, stagnation, nonfinite values, or memory limits explicitly.

\subsection{Assembly reuse and qualification}
The useful assembly interface is the oriented face-BSR graph, dense local
couplings, and reusable local factorizations. Factor $L_K$ once, apply it
to multiple trace/source columns, and optionally retain factors or recovery
maps for reconstructing $\underline{q}^K$ and $\underline{u}^K$ after solving
for $\underline{\widehat{u}}$. Reuse requires an unchanged local
operator, including coefficients, stabilization, geometry, and time-step
terms. Weigh retained memory against recomputation.

An assembly kernel may insert the modal trace transform directly in local
couplings before scattering to precomputed BSR slots. Element-face incidence
maps can also provide initial patch neighborhoods. These changes save setup
or recovery work without needing a PDE-regime detector.

Invertible changes confined to the eliminated interior variables leave the
exact Schur complement unchanged. For example, with $y=X\widetilde y$ and
test transform $Y$, the transformed interior data are $YL_KX$, $YB_K$,
and $C_KX$, so
\[
 D_K-C_KX(YL_KX)^{-1}YB_K=D_K-C_KL_K^{-1}B_K.
\]
Local equilibration can therefore improve finite-precision factorization,
but does not itself improve the exact trace spectrum. Stabilization changes
the discretization and should not be an implicit preconditioning adjustment.

Keep the current SPD path available. A separate nonsymmetric hierarchy can
reuse modal/BSR primitives and patch extraction, while adding independent
transfer storage, block sparse triple products, nonsymmetric local factors,
a nonsymmetric AMG endpoint, and flexible right-preconditioned Arnoldi.
The symmetric-part adapter in \code{linalg/multigrid/krylov.py} is a useful comparator,
not this algorithm.

The accompanying small-matrix diagnostic verifies the ideal identities,
the singular-truncation counterexample, rejection of a singular high-mode
split, and fixed-cycle right-preconditioned solves on four synthetic block
matrices. It uses dense exact terminal solves and fixed block-Jacobi high
relaxation. Thus it verifies algebra only: it tests neither AIR below the
p levels, the proposed inner-GMRES smoother, GPU performance, nor HDG
mesh/degree robustness. Results and commands are preserved in
\code{docs/research/solver_studies/nonsymmetric_pmg_design_2026_09_22.md}.

Production qualification must separately measure original-system residuals,
setup and solve time, iteration growth under h/p refinement, coarse/transfer
storage, and peak memory. Include nonsymmetric systems with difficult
symmetric parts, unequal couplings, cyclic graphs, face orientation changes,
and coefficient scaling. Use one algebraic policy across cases before
claiming generality. If polynomial coarsening repeatedly fails these checks,
starting block AMG earlier is preferable to assuming every problem must
reach a scalar p=0 stage.

\section{Poisson h/p hierarchy alternatives}\label{sec:hpchoices}
The geometric candidate retains linear traces after p-coarsening and then
coarsens the skeleton by cell agglomeration. The algebraic candidate retains
these same degree-one traces before face aggregation and interpolation
optimization. Both use the HDG operator and coordinates from
Section~\ref{sec:shared}. They extend the hierarchy choices, not the supported
API of Section~\ref{sec:spd}. No new HDG convergence or performance measurements
are claimed for these candidates.

The closest direct HDG template is Wildey, Muralikrishnan, and Bui-Thanh
\cite{wildey}: p-to-one reduction, agglomeration, DtN-based transfers and
local interior correction. The initial HYBRIDGE comparison should use
$p\to1$ against a staged schedule such as $6\to3\to1$. Same-mesh p-transfer
can reuse the Galerkin injection from Section~\ref{sec:spd}, terminating at
one. Each matrix inherits the original condensed operator and stabilization;
this does not rediscretize a lower-order HDG problem.

\subsection{Harmonic coarsening and macro-edge spaces}\label{sec:harmonic}
At one h step, split trace coefficients into interior faces $I$ of the next
macro-elements and their retained boundary faces $B$. Let $J_\ell$ map
coarse macro-boundary coefficients to the retained fine boundary. For SPD
$A_\ell$, set
\[
 A_\ell=\begin{bmatrix}A_{II}&A_{IB}\\A_{BI}&A_{BB}\end{bmatrix},\qquad
 P_\ell=\begin{bmatrix}-A_{II}^{-1}A_{IB}J_\ell\\J_\ell\end{bmatrix},
 \qquad R_\ell=P_\ell^T.
\]
Then
\[
 A_{\ell+1}=J_\ell^T
 (A_{BB}-A_{BI}A_{II}^{-1}A_{IB})J_\ell,\qquad
 C_{\ell,I}=\begin{bmatrix}A_{II}^{-1}&0\\0&0\end{bmatrix}.
\]
These are the SPD assembled-matrix specialization of the DtN construction
\cite{wildey}. The extension satisfies $(A_\ell P_\ell)_I=0$ and minimizes
fine energy for prescribed retained-boundary values. The Galerkin energy
identity follows: $(P_\ell v)^TA_\ell(P_\ell v)=v^TA_{\ell+1}v$.
The interior block is local to each aggregate when the inherited trace
coupling follows cell-local assembly. Verify that locality at every level;
it is not guaranteed for an unrelated AMG coarse matrix.

The local correction is essential to the proposed starting cycle. Harmonic
prolongation alone does not recover arbitrary interior error, and some
block-Jacobi configurations fail without the correction in \cite{wildey}.
A symmetric h-level application starts from zero, pre-smooths, applies
$C_{\ell,I}$ to the residual, applies the recursive coarse correction,
applies $C_{\ell,I}$ again, and post-smooths with the adjoint sequence.
With exact factors, $P_\ell^TA_\ell C_{\ell,I}=0$; the two placements of the
interior correction do not double its exact action. Inexact local solves
require rechecking the symmetry and positivity of the complete cycle.
Use the existing p-level cycle without this interior correction on p-only
transitions.

For a first 2D construction, group connected cells into aggregates of about
four current cells, keeping diameter and interface complexity controlled.
This is a proposed starting policy, not an optimized published prescription.
Split disconnected or branched interfaces and respect holes and boundary
tags. Represent each macro-edge $E$ by two reference-normalized modes in a
fixed arclength parameter. On child edge $f$, set
\[
 J_f=M_f^{-1}C_{fE},\qquad
 (C_{fE})_{ij}=\int_f\psi_i^f\psi_j^E\,ds.
\]
Linear arclength traces restrict exactly to linear functions on straight
child segments. Include physical length factors and the sign of odd modes
under coordinate reversal.

Two arclength modes on a bent macro-edge cannot generally represent both
physical $x$ and $y$ traces. Prefer splitting strongly bent interfaces;
enrichment by independent traces of $\{1,x,y\}$ is another option but can
produce nonuniform block ranks. Lu, Rupp, and Kanschat \cite{luinjection}
provide useful stability and conforming-linear-trace reproduction criteria,
as well as interpolation and local-reconstruction injections on refined
meshes. Their h-convergence analysis does not establish these properties for
arbitrary agglomerates or prove p-uniform convergence of this proposal.

\subsection{Smoothing, storage, and implementation choices}
For the existing GPU setting, the proposed first smoother combines full
face-block diagonal solves with degree-two Chebyshev acceleration and
balanced $1+1$ sweeps. Compare higher degree or stronger patches by complete
cycle cost. This combination is an HYBRIDGE adaptation; it is not asserted to
be the optimum from \cite{wildey}. The spectral cautions in
Section~\ref{sec:spd} still apply. Vertex-patch and harmonic-transfer ideas
also appear in the divergence-conforming HDG work of Fu and Kuang
\cite{fuhp}, but its flow spaces and kernel argument differ from Poisson.

The direct degree-six schedule has block dimensions $7\to2\to2\to\cdots$;
a staged schedule has $7\to4\to2\to2\to\cdots$. Only the h steps reduce
face count. Choose the number of levels from actual coarse dimensions and
factorization cost; four-cell aggregation does not ensure a fourfold trace
reduction. An initial terminal-size comparison is 128--512 scalar unknowns
with a cached direct factorization. This is a tuning range, not a theorem.

Polynomial transfers can remain implicit mode copies. Harmonic h transfers
are globally rectangular, even with uniform $2\times2$ blocks. Their
application may store $P_\ell$ or use $J_\ell$, $A_{IB}$ and aggregate factors;
restriction must use the corresponding transpose action. The current native
p-cycle assumes a fixed face count, so this requires a generalized level
interface. Schur elimination creates boundary coupling within aggregates:
count scalar storage, transfer fill, local factors and temporary workspaces,
not only block entries. Variable degree and 3D macro-faces require separate
choices.

\subsection{Alternative coarse spaces and comparison criteria}
The algebraic alternative uses $p\to1$ followed by whole-face aggregation
and constrained energy-minimizing interpolation. Dobrev et al.
\cite{dobrev} motivate AMG on hybridized systems; Olson and Schroder
\cite{olson} provide block-relaxation and interpolation-optimization
precedent for volume DG Poisson. Their results do not qualify this particular
HDG trace adaptation.

In the reference-normalized coordinates, the physical constant trace has
coefficients $(\sqrt{2},0,\ldots)$ on each face, up to global normalization.
It is not the all-ones modal vector. Dirichlet elimination means it is an
interior smooth-error candidate rather than an exact global null vector.
Respect homogeneous error boundary data. Rank-test aggregate restrictions of
candidate vectors, form tentative interpolation, and minimize column energy
on a bounded sparsity pattern subject to $PN_c=N_f$. Physical coordinate
traces or independently relaxed errors are possible additional candidates.
Use $R=P^T$ and Galerkin operators for the SPD proposal. Algebraic coarse
blocks count independent candidates, not necessarily polynomial modes;
variable candidate ranks need grouped or variable-block storage.

A further candidate uses the trace of a conforming P1 auxiliary space,
with its Galerkin matrix $P^TA_1P$ and scalar AMG. Hybridized mixed multigrid
\cite{gopal} is relevant precedent; substituting a separately assembled
P1 stiffness requires an equivalence argument. Fu and Kuang's HDG-P0/CR
result \cite{fupzero} gives a different valid lowest-order construction.
The high-order operator's principal constant-mode block is not automatically
that special quadrature-based HDG-P0 operator. Enriched multilevel trace
spaces \cite{muralilevel} provide another direction when small coarse spaces
miss significant error.

The literature review and retained-study links are in
\code{docs/algorithms/hp_amg/hierarchies.md}. Current scalar-AMG repacking
and wider finest-patch studies motivate inspecting coarse-space quality
before adding storage or smoothing cost. They do not reject every block
hierarchy or patch construction.

A later comparison must hold the fine operator, boundary conditions,
stabilization, load, initial guess and true-residual target fixed. Measure
setup, repeated solve time, iterations under h/p refinement, and hierarchy
memory. Small assembled HDG systems must check orientation, transfer rank,
constant and applicable affine reproduction, Galerkin identities, complete
cycle symmetry/positivity, and nullspaces. Pure Neumann or periodic Poisson
needs compatible loads and a consistent gauge or quotient-space treatment.
No new implementation, HDG benchmark, build, or compilation is part of this
documentation unification.

\begin{thebibliography}{99}
\bibitem{lair} B.~S. Southworth, T.~A. Manteuffel, and J. Ruge,
\emph{Nonsymmetric Algebraic Multigrid Based on Local Approximate Ideal
Restriction}, SIAM J. Sci. Comput. 40 (2018), A4105--A4130.
\url{https://arxiv.org/abs/1708.06065}.
\bibitem{hdgair} A.~A. Sivas, B.~S. Southworth, and S. Rhebergen,
\emph{AIR algebraic multigrid for a space-time hybridizable discontinuous
Galerkin discretization of advection(-diffusion)}.
\url{https://arxiv.org/abs/2010.11130}.
\bibitem{clair} A. Ali et al.,
\emph{Constrained Local Approximate Ideal Restriction for Advection-Diffusion
Problems}, SIAM J. Sci. Comput. (2024).
\url{https://arxiv.org/abs/2307.00229}.
\bibitem{fgmres} PETSc, \emph{KSPFGMRES} documentation.
\url{https://petsc.org/release/manualpages/KSP/KSPFGMRES/}.
\bibitem{pyamg} PyAMG, \emph{Classical AMG: air\_solver} documentation.
\url{https://pyamg.readthedocs.io/en/latest/generated/pyamg.classical.html}.
\bibitem{hypre} hypre, \emph{ParCSR Solvers:
HYPRE\_BoomerAMGSetRestriction} API documentation.
\url{https://hypre.readthedocs.io/en/latest/api-sol-parcsr.html}.
\bibitem{wildey} T. Wildey, S. Muralikrishnan, and T. Bui-Thanh,
\emph{Unified Geometric Multigrid Algorithm for Hybridized High-Order
Finite Element Methods}, SIAM J. Sci. Comput. 41 (2019), S172--S195.
\url{https://arxiv.org/html/1811.09909}.
\bibitem{luinjection} P. Lu, A. Rupp, and G. Kanschat,
\emph{Analysis of Injection Operators in Geometric Multigrid Solvers for
HDG Methods}, SIAM J. Numer. Anal. 60 (2022), 2293--2317.
\url{https://arxiv.org/pdf/2104.00118}.
\bibitem{gopal} J. Gopalakrishnan and S. Tan,
\emph{A Convergent Multigrid Cycle for the Hybridized Mixed Method},
Numer. Linear Algebra Appl. 16 (2009), 689--714.
\url{https://web.pdx.edu/~gjay/pub/hmg.pdf}.
\bibitem{fupzero} G. Fu and W. Kuang,
\emph{Optimal Geometric Multigrid Preconditioners for HDG-P0 Schemes
for the Reaction-Diffusion Equation and the Generalized Stokes Equations},
ESAIM M2AN 57 (2023), 1553--1587.
\url{https://arxiv.org/abs/2208.14418}.
\bibitem{fuhp} G. Fu and W. Kuang,
\emph{hp-Multigrid Preconditioner for a Divergence-Conforming HDG Scheme
for the Incompressible Flow Problems}, J. Sci. Comput. (2024).
\url{https://doi.org/10.1007/s10915-024-02568-4}.
\bibitem{dobrev} V. Dobrev, T. Kolev, C. S. Lee, V. Tomov, and P. S. Vassilevski,
\emph{Algebraic Hybridization and Static Condensation with Application
to Scalable H(div) Preconditioning}, SIAM J. Sci. Comput. (2019).
\url{https://arxiv.org/html/1801.08914}.
\bibitem{olson} L. N. Olson and J. B. Schroder,
\emph{Smoothed Aggregation Multigrid Solvers for High-Order Discontinuous
Galerkin Methods for Elliptic Problems}, J. Comput. Phys. 230 (2011), 6959--6976.
\url{https://lukeo.cs.illinois.edu/files/2011_OlSc_hodg.pdf}.
\bibitem{muralilevel} S. Muralikrishnan, T. Bui-Thanh, and J. N. Shadid,
\emph{A Multilevel Approach for Trace System in HDG Discretizations},
J. Comput. Phys. 407 (2020), 109240.
\url{https://arxiv.org/abs/1903.11045}.
\end{thebibliography}
\end{document}
